%% notes-tex/025-additive-baselines/sections/7-remaining.tex

\section{What is established, and what is left\secstatus{todo}}
\label{sec:remaining}

\subsection{Established\secstatus{todo}}

The additive null is fit on the same records, splits and label the transformer
uses, with the ridge penalty chosen on validation and reported on test, under both
arms. On the published split the transformer's margin over a model that cannot
represent interaction is real and small, and most of what every model scores there
is screen identity: a model that never looks at the third gene reaches 97 percent
of the additive null, and holding out whole query doubles drops the additive null
from 0.400 to 0.185 on arm Q and to $0.127 \pm 0.033$ over five folds, and inverts
the ordering of the baselines.\iflong{} Every model in the comparison receives the
same information, the identity of the three perturbed genes; that is a reading of
the code, verified by a hand-built forward pass that reproduces the recorded
metrics. The graphs are load-bearing during training: they carried 99.99 percent of
the 010 loss, their heads recover their graphs, and the matched pair of runs without
the penalty does not train.\fi

\subsection{Not established\secstatus{todo}}

\notesec{The transformer's capability on held-out screens is not measured}
One checkpoint of one run has a test score on arm Q, and it is below the additive
null with an interval that excludes zero (Table~\ref{tab:decision}). That is a
measurement of GH 1640's epoch 7 checkpoint. Whether the configuration can do
better, or whether a seed-averaged run lands where this one did, is D3, and it has
not been run. Whether the $+0.04$ margin of the random split survives held-out
screens is therefore answered for one draw of the optimizer and open for the
model.

\iflong
\notesec{The reproduction is faithful but not bit-exact}
The hand-built forward pass matches the recorded 010 metrics to the fourth decimal,
not exactly. The original runs used mixed bf16 precision and the reproduction runs
in fp32, which is the likely source and has not been verified by rerunning in
bf16. The conclusions of Section~\ref{sec:paired} rest on correlations at the
second decimal and are not sensitive to a fourth-decimal shift.

\notesec{One seed of the nonlinear baseline enters the paired analysis}
Table~\ref{tab:paired} uses B5 seed 0. The three seeds agree closely on aggregate
metrics, but their record-level agreement with each other was not measured.

\notesec{Whether the graphs carry information is separate from whether they are
load-bearing}
The runs without the graph penalty fail to train, and the penalty was almost the
entire loss. Neither separates the graphs supplying biological structure from a
large auxiliary term supplying conditioning that the optimization needs. The
control that would separate them, a comparable penalty against degree-matched
random graphs, has not been run, and the matched pair itself ran 30 epochs with
one diverged arm rather than three replicates.

\notesec{Nothing here bears on the architecture's merits elsewhere}
The set-function reduction of Section~\ref{sec:cgt} is specific to the 010
configuration, where the masking and propagation arms are off and the
perturbation transform is one block. A configuration that switches those on would
have a strain-dependent encoder path and would need its own reading.
\fi

\subsection{What is left before this enters the main text\secstatus{todo}}

\begin{enumerate}
\item \textbf{Replicate the 010 configuration on arm Q.} Three seeds of
  \texttt{cgt\_s0\_q\_kl\_004} with a test evaluation at the best validation epoch,
  under one fixed epoch budget. GH 1640 peaked at epoch 7 and 36 epochs cost the 12
  hour wall clock, so a 20 epoch budget with checkpoint-on-best is enough. These
  fill the empty slots of Figure~\ref{fig:disjointruns}c, decide D3 of
  Table~\ref{tab:decision}, and give the transformer row of
  Table~\ref{tab:heldout} a spread. The CPU scorer and the paired bootstrap run
  unchanged on each.
\item \textbf{Run the missing control.} \texttt{cgt\_s0\_q\_kl\_ctrl\_016}, the
  learnable table under the IGB schedule and readout, with seeds and test scores;
  without it the embedding variants of Table~\ref{tab:wandbruns} stay validation
  trajectories and do not enter Figure~\ref{fig:ladders}.
\item \textbf{Fold spread for the transformer, or an explicit statement that it is
  unaffordable.} Panel d of Figure~\ref{fig:ladders} is five folds for the nulls. A
  transformer fold costs about 12 hours of one GilaHyper node per 36 epochs, so five
  folds by three seeds is 15 such runs. If the main text reports arm Q alone it
  should say so and cite the fold spread of the null as the scale of the
  split-to-split variation.
\item \textbf{B5 on GPU.} The 025 B5 fit ran on CPU and differs from 010's GPU fit
  by 0.006 on arm R. A GPU refit closes the one row of Table~\ref{tab:heldout} that
  is not an exact reproduction.
\iflong
\item \textbf{The graph control.} A run with the penalty of Eq.~\ref{eq:graphreg}
  against degree-matched random graphs, to separate the graphs carrying structure
  from a large penalty carrying conditioning.
\fi
\item \textbf{Bibliography.} The Visani preprint, Tran 2026 and Kuzmin 2018 are
  named by DOI in prose. A Zotero collection for this document, then
  \texttt{make bib-pull}, makes them citations.
\end{enumerate}

\iflong
\notesec{What is not carried over from the 010 report}
The gene-index defect of the 010 inference runs, a uniform 28-position shift that
invalidated the design-space selections, is in \texttt{notes-tex/010-index-defect}
and the 010 report's design-space section, and bears on the panel selection rather
than on the ladder.
\fi
